\chapter{Termination}
\index{Termination}

\medskip
\noindent\textit{Educational reference; always seek professional advice for individual care.}

\section*{Definition and Overview}
Termination of pregnancy, commonly called abortion, is the medical or surgical ending of a pregnancy before birth. It is a routine part of reproductive health care in many countries, where it is lawful in all states and territories, although the gestational limits and administrative arrangements differ between jurisdictions. Termination may be provided using medication alone, known as medical termination, or by a surgical procedure, and in some cases a combination of both. When delivered by trained professionals in an appropriate clinical setting, termination is extremely safe, and serious complications are rare. This chapter explains the available methods, what happens before, during and after the procedure, the warning signs that require urgent medical attention, and the longer-term outlook. It is intended for anyone considering or supporting someone through a termination, and aims to present balanced, non-judgemental information that supports informed decision-making.

\section*{Epidemiology}
Termination of pregnancy is one of the most common gynaecological procedures worldwide. It is estimated that roughly one in four women will have a termination at some point during their reproductive years, a figure broadly consistent with other high-income countries. The majority of terminations are performed in the first trimester, within the first twelve weeks of pregnancy, and medical termination now accounts for a substantial and growing proportion of procedures in many settings. Women seek termination for many different reasons, including contraceptive failure, concerns about maternal health, the diagnosis of a serious fetal abnormality, and personal, social, financial or relationship circumstances. Access is generally easier in metropolitan areas than in rural and remote parts of some countries, and cost, distance, and limited provider availability remain real barriers for some women. Because unsafe abortion is a leading cause of maternal death and disability worldwide, access to safe, legal termination services is recognised by the World Health Organization as a fundamental component of sexual and reproductive health care.

\section*{Aetiology and Pathophysiology}
Termination is not a disease, and it has no single "cause". Rather, it is a response to a pregnancy that a person does not wish to continue, whatever the reason. Understanding why terminations happen is helpful because it frames the care that should be offered.

\begin{itemize}
    \item \textbf{Unintended pregnancy:} the most common scenario, often due to contraceptive failure such as missed pills, a broken condom, or a late injection, or due to using no contraception at all.
    \item \textbf{Maternal health reasons:} continuing a pregnancy may carry significant risk in the presence of severe cardiac, renal, psychiatric, or other medical conditions, and termination may be recommended to protect the pregnant person.
    \item \textbf{Fetal abnormality:} a serious structural or chromosomal abnormality may be diagnosed on early screening or diagnostic testing, prompting a difficult decision about whether to continue.
    \item \textbf{Personal and social circumstances:} financial pressures, unstable housing, relationship breakdown, caring responsibilities, or being at a particular stage of life or study.
    \item \textbf{Sexual violence:} pregnancy resulting from rape or coercion, in which the decision to terminate is made in a context of significant trauma and requires sensitive, specialised support.
\end{itemize}

The decision to end a pregnancy is deeply personal and emotionally complex. Whatever the reason, a person presenting for termination is entitled to respectful, confidential, and non-judgemental care, and to counselling support if they want it.

\section*{Clinical Features}
The experience of a termination varies according to the method used and the stage of pregnancy, but there are a number of common features that anyone having the procedure should expect.

\begin{itemize}
    \item \textbf{Medical termination} involves taking medication that ends the pregnancy, followed by cramping and bleeding that typically begins within hours and continues over several days, often with passage of the pregnancy tissue.
    \item \textbf{Surgical termination} is a brief procedure that removes the pregnancy, usually under local anaesthesia with sedation or a general anaesthetic, followed by mild cramping and light to moderate bleeding.
    \item \textbf{Bleeding} after either method is usually similar to a heavy period, settles within one to two weeks, and gradually becomes lighter.
    \item \textbf{Pain} is normally controlled with simple analgesia such as paracetamol and ibuprofen, with most people rating the cramping as similar to period pain.
    \item \textbf{Nausea, breast tenderness and fatigue} may persist briefly as pregnancy hormones fall after the termination.
    \item \textbf{Emotional responses} vary widely. Most people feel relief, but sadness, guilt, or mixed emotions are also normal and usually settle with time and support.
    \item \textbf{Return to normal activities} is generally rapid, with most people able to return to work, study, or home duties within one to three days.
\end{itemize}

\section*{Differential Diagnosis}
Before a termination is provided, the situation must be confirmed and certain alternative diagnoses excluded, because each requires a different pathway of care.

\begin{itemize}
    \item \textbf{Confirmed intrauterine pregnancy:} the diagnosis to be established before any termination, including accurate dating of gestational age.
    \item \textbf{Miscarriage:} a pregnancy that has already ended spontaneously; a person may present for termination only to discover the pregnancy has miscarried, in which case management changes.
    \item \textbf{Ectopic pregnancy:} a pregnancy implanted outside the womb, usually in the fallopian tube. This is a medical emergency that requires entirely different, urgent treatment, and a termination procedure will not resolve it.
    \item \textbf{Molar pregnancy:} an abnormal pregnancy in which trophoblastic tissue grows abnormally; it requires specialist monitoring and follow-up rather than standard termination care.
    \item \textbf{Ongoing pregnancy after attempted termination:} more likely when medication has not been taken as directed or when the pregnancy is later than thought, and is confirmed by follow-up ultrasound.
    \item \textbf{Incomplete termination:} retained pregnancy tissue after a procedure, which may cause continuing bleeding and requires review and sometimes further treatment.
\end{itemize}

\section*{When to Seek Medical Care}
After a termination, the clinic or a general practitioner should be contacted promptly for heavy bleeding, severe pain, fever, or anything that suggests infection or an incomplete procedure. A follow-up review is routinely arranged, usually two to three weeks later, to confirm that the termination is complete and that recovery is normal.

\textbf{Contact your local emergency services for:}
\begin{itemize}
    \item Very heavy bleeding that soaks more than two sanitary pads in an hour or passes large clots
    \item Severe, worsening, or one-sided abdominal or pelvic pain
    \item Fever, chills, or foul-smelling vaginal discharge
    \item Fainting, light-headedness, or collapse
    \item Chest pain or difficulty breathing
\end{itemize}

These symptoms can indicate haemorrhage, infection, or complications of the procedure, and need urgent assessment.

\section*{Diagnosis and Investigations}
Assessment before a termination is straightforward and is designed to confirm the pregnancy, establish the gestational age, and ensure that the chosen method is safe.

\begin{itemize}
    \item \textbf{Pregnancy test:} a urine or blood test to confirm the pregnancy.
    \item \textbf{Ultrasound:} to confirm that the pregnancy is inside the womb rather than ectopic, to date it accurately, and to guide the choice of method.
    \item \textbf{Blood tests:} full blood count, and blood group testing so that Rh-negative people can receive anti-D immunoglobulin to prevent sensitisation that could harm future pregnancies.
    \item \textbf{Gestational age assessment:} the method offered depends on how many weeks pregnant the person is, with gestational limits set by local law and clinical guidelines.
    \item \textbf{Screening for infection:} testing for chlamydia and gonorrhoea is often recommended, because untreated infection can spread during the procedure.
    \item \textbf{Counselling and consent:} a full discussion of methods, risks, benefits, and after-care, with time to consider the decision.
\end{itemize}

\section*{Management}
The method of termination is chosen according to gestational age, availability, medical history, and the person's preference.

\begin{itemize}
    \item \textbf{Medical termination:} mifepristone followed by misoprostol, generally used in the first trimester up to about nine weeks, and taken at home or in the clinic depending on local arrangements.
    \item \textbf{Surgical termination:} suction aspiration, usually performed in the first trimester under local anaesthesia with sedation or general anaesthetic; it is quick, safe, and associated with a low rate of incomplete termination.
    \item \textbf{Second-trimester procedures:} for pregnancies beyond fourteen weeks, surgical management is usually required, sometimes performed over one or two days, and may involve induction of labour with medication.
    \item \textbf{Pain relief and comfort measures:} simple analgesia, heat packs, and emotional support during and after the procedure.
    \item \textbf{Contraception:} advice and, where requested, provision of an effective method so that the person can start contraception immediately, as fertility returns quickly.
    \item \textbf{Psychological support:} counselling before and after, including referral to specialised services for those who need it.
    \item \textbf{Follow-up:} review to confirm completion and normal recovery, and to address any physical or emotional concerns.
\end{itemize}

\section*{Complications}
Serious complications are uncommon when a termination is performed safely, but they can occur and it is important to recognise them.

\begin{itemize}
    \item Bleeding, including incomplete termination with retained tissue, which may require a further procedure
    \item Infection of the uterus or pelvis, which usually responds to antibiotics
    \item Injury to the cervix or, rarely, perforation of the uterus during a surgical termination
    \item Ongoing pregnancy when a medical termination does not complete, detected at follow-up
    \item Asherman's syndrome, scarring of the uterine lining after surgery, which is uncommon but can affect future fertility and menstruation
    \item Emotional responses, including sadness, guilt, or regret, in a minority of people
    \item Complications are far more common with unsafe or unregulated procedures, which is why access to safe, legal care matters so much
\end{itemize}

\section*{Prognosis and Outcomes}
Safe termination is one of the lowest-risk medical procedures, and the great majority of people recover fully with no lasting physical effects. Fertility is preserved, and most people who subsequently wish to conceive go on to have normal pregnancies. Physical recovery is usually rapid, with bleeding settling within two weeks and normal activities resuming quickly. Research consistently shows that most people do not regret their decision, and the very small proportion who experience distress generally improve with support and time. There is no evidence that a single termination causes long-term harm to physical or mental health, and future fertility is not affected by an uncomplicated first-trimester termination.

\section*{Prevention and Self-Care}
\begin{itemize}
    \item Use effective contraception consistently and correctly if pregnancy is not desired
    \item Be familiar with emergency contraception for use within the time limits after unprotected sex
    \item Attend the planned follow-up appointment so that any problem is detected early
    \item Rest and use simple analgesia for cramping in the days after the procedure
    \item Avoid tampons, swimming, and penetrative sex until bleeding has settled, to reduce infection risk
    \item Arrange contraception immediately after a termination, as ovulation can resume within weeks
    \item Seek counselling or support if feelings are mixed or difficult, and be aware that this is normal
    \item Know the emergency warning signs and where to get help, day or night
\end{itemize}

\section*{Sources and Further Reading}
The following organisations provide reliable, non-judgemental information about termination of pregnancy for people making this decision, their families, and clinicians.

\begin{itemize}
    \item NHS. \textit{Abortion (termination of pregnancy): what to expect.} https://www.nhs.uk/conditions/abortion/
    \item World Health Organization. \textit{Abortion care and safe abortion guidance.} https://www.who.int/
    \item MSD Manual. \textit{Overview of abortion: methods and complications.} https://www.msdmanuals.com/
    \item MedlinePlus. \textit{Abortion health information.} https://medlineplus.gov/abortion.html
    \item Family Planning Australia. \textit{Reproductive and sexual health information.} https://www.fpnsw.org.au/
    \item Children by Choice. \textit{Information and counselling for unplanned pregnancy.} https://www.childrenbychoice.org.au/
\end{itemize}
